\documentclass[10pt]{article}

\usepackage[utf8]{inputenc}

\usepackage[T1]{fontenc}

\usepackage{amsmath}

\usepackage{amsfonts}

\usepackage{amssymb}

\usepackage[version=4]{mhchem}

\usepackage{stmaryrd}



\DeclareUnicodeCharacter{0131}{$\imath$}



\begin{document}

Во многих случаях задание Р. вероятностей указанием возможных значений случайной величины и соответствующих им вероятностей невыполнимо. Напр., если случайная величина принимает любые значения из отрезка $[-1 / 2,1 / 2]$, го Р. вероятностей задается указанием вероятностей того, что случайная величина примет значение из любого заданного интервала, принадлежащего отрезку $[-1 / 2,1 / 2]$, так как вероятность каждого отдельного значения может быть равна нулю. Если существует такая функция $p(x)$, что вероятность попадания $X$ в любой интервал $(a, b)$ на прямой равна



\[

\mathrm{P}\{a<X<b\}=\int_{a}^{b} p(x) d x,

\]



причем $p(x) \geqslant 0$ и



\[

\int_{-\infty}^{+\infty} p(x) d x=1

\]



го Р. вероятностей величины $X$ называется а бс олютно н е п р е р ы в ны м, а функция $p(x)$ носит название плотности вероятности.



В случае равномерного распределения на отрезке $[-1 / 2,1 / 2]$



\[

p(x)= \begin{cases}1, & |x| \leqslant 1 / 2 \\ 0, & |x|>1 / 2\end{cases}

\]



Важнейшее Р. вероятностей не преры вного ти паГаусса распределение с плотностью



\[

p(x)=\frac{1}{\sqrt{2 \pi} \sigma} e^{-(x-a)^{2} / 2 \sigma^{2}},

\]



де $-\infty<a<\infty$ и $\sigma>0$. Р. вероятностей не исчерпываются ıискретным и непрерывным типами: они могут быть и јолее сложной природы. Описание Р. вероятностей, приодное во всех случаях, может быть, напр., достигнуто при Іомощи распределения функции, к-рая для любой случайной зеличины $X$ при каждом действительном $x$ определяется рормулой



\[

F(x)=\mathrm{P}\{X<x\} .

\]



Вероятность того, что $X$ примет значение из нек-рого ингервала $[a, b)$ на прямой, равна



\[

\mathrm{P}\{a \leqslant X<b\}=F(b)-F(a) .

\]



Функция Р. однозначно задает соответствующее Р. вероятюостей случайной величины.



Часто полное описание Р. вероятностей (напр., при поиощи плотности или функции Р.) заменяют заданием нејольшого числа числовых характеристик, к-рые указывают, қак правило, наиболее типичные (в том или ином отношеıии) значения случайной величины и степень рассеяния іначений случайной величины около нек-рого типичного เначения. Из этих характеристик наиболее употребительны иатематическое ожидание (среднее значение), дисперсия и иомент.



Помимо функций Р. полную информацию о Р. вероятнотей содержат так наз. производящие функции и характери'тические функции. Производящая функция Р. вероятно:тей случайной величины $X$ с целыми неотрицательными іначениями определяется формулой



\[

\varphi(t)=\mathrm{E} z^{X}

\]



Іри каждом фиксированном $z,|z| \leqslant 1$. Характеристич. фунщия Р. вероятностей произвольной случайной величины $X$ пределяется по формуле



\[

f(t)=\mathrm{E} e^{i t X}

\]



ля любого действительного $t,-\infty<t<\infty$. Производящие и арактеристич. функции однозначно определяют соответстเуюшие Р. вероятностей.



Если $X_{1}, \ldots, X_{n}$ - несколько случайных величин, опредегенных на одном и том же вероятностном пространстве, о говорят о совместном Р. вероятностей величин $\mathrm{Y}_{1}, \ldots, X_{n}$ или о $\mathrm{P}$. вероятностей случайного вектора $\mathrm{Y}=\left(X_{1}, \ldots, X_{n}\right)-$ многоме рном Р. вероятностей. Мно- гомерное Р. вероятностей задается набором вероятностей вида:



\[

\mathrm{P}\left\{X_{1}=x_{1}, \ldots, X_{n}=x_{n}\right\} \text { (дискретный тип) }

\]



или с помощью плотности



\[

p\left(x_{1}, \ldots, x_{n}\right) \text { (непрерывный тип) }

\]



или же в общем случае функцией Р. Если случайные величины $X_{1}, \ldots, X_{n}$ независимы, то Р. вероятностей $X=\left(X_{1}, \ldots\right.$, $X_{n}$ ) определяется Р. вероятностей отдельных величин $X_{k}$, $k=1, \ldots, n$. В ином случае приходится рассматривать усл о в н ы е Р. вероятностей одних величин при фиксированных значениях других и соответствующие числовые характеристики (условные математич. ожидания, условные дисперсии, смешанные моменты). См. также Корреляция.



Статистич. аналогом P. вероятностей служит так наз. эмпирическое распределение. Пусть $X_{1}, \ldots, X_{n}-$ результаты наблюдений, к-рые предполагаются независимыми, одинаково распределенными случайными величинами с функцией Р. вероятностей $F(x)$. Эмпирич. Р. вероятностей характеризуется так наз. эм п и ричес кой функ ц и е й распределения



\[

F_{n}^{*}(x)=n_{x} / n,

\]



где $n_{x}$ - число наблюдений с $X_{k}<x$, к-рой соответствует так наз. выборочное среднее



\[

\bar{X}=\frac{1}{n} \sum_{k=1}^{n} X_{k}

\]



и выборочная дисперсия



\[

s^{2}=\frac{1}{n} \sum_{k=1}^{n}\left(X_{k}-\bar{X}\right)^{2}

\]



Эмпирич. Р. и его характеристики используются для приближенного представления теоретич. Р. вероятностей и его характеристик.



Лит.: [1] Феллер В., Введение в теорию вероятностей и ее приложения, пер. с англ., 3 изд., т. 1-2, М., 1984; [2] П рохоро в Ю. В., Р о за н о в Ю.А., Теория вероятностей, 3 изд., М., 1987; [3] Справочник по теории вероятностей и математической статистике, 2 изд., М., $1985 . \quad$ А. В. Прохоров. РАСПРЕДЕЛЕНИЯ ВЕКТОР к в а то в Й - аналог классического вектора распределения (см. Распределения вектор к л а с с и ч е с к и й) в квантовой статистической механике. Квантовым вектором распределения системы $N$ тождественных ферми- или бозе-частиц называется совокупность $s$-частичных функций Вигнера:



\[

\begin{gathered}

f=\left\{f_{s}^{B}\left(\overrightarrow{q_{1}}, \overrightarrow{p_{1}} ; \ldots ; \overrightarrow{q_{s}}, \overrightarrow{p_{s}}\right), s=0,1,2, \ldots, N\right\}, \\

f_{s}^{B}\left(\overrightarrow{q_{1}}, \overrightarrow{p_{1}} ; \ldots ; \vec{q}_{s}, \overrightarrow{p_{s}}\right)=\left(8 \pi^{3}\right)^{-s} \int d \overrightarrow{k_{1}} \ldots d \overrightarrow{k_{s}} \times \\

\times \exp \left(-i \sum_{n=1}^{s} \overrightarrow{k_{n}} \overrightarrow{p_{n}}\right) \operatorname{Tr}\left(\rho \Psi^{+}\left(\overrightarrow{q_{1}}-\frac{\hbar}{2} \overrightarrow{k_{1}}\right) \ldots \Psi^{+}\left(\overrightarrow{q_{s}}-\frac{\hbar}{2} \overrightarrow{k_{s}}\right) \times\right. \\

\left.\times \Psi\left(\overrightarrow{q_{s}}+\frac{\hbar}{2} \overrightarrow{k_{s}}\right) \ldots \Psi\left(\overrightarrow{q_{1}}+\frac{\hbar}{2} \overrightarrow{k_{1}}\right)\right),

\end{gathered}

\]



где $\Psi^{+}, \Psi$ - соответствующие операторные функции, фигурирующие в формализме вторичного квантования, $\rho$ - статистич. оператор системы. Компоненты квантового Р. в. имеют нормировку:



\[

\int d \overrightarrow{q_{1}} d \overrightarrow{p_{1}} \ldots d \overrightarrow{q_{s}} d \overrightarrow{p_{s}} f_{s}^{B}\left(\overrightarrow{q_{1}}, \overrightarrow{p_{1}} ; \ldots ; \overrightarrow{q_{s}}, \overrightarrow{p_{s}}\right)=\frac{N !}{(N-s) !},

\]



совпадающую с нормировкой $s$-частичных классич. приведенных распределения функций. Эволюция во времени квантового Р. в. определяется квантовым обобщением цепочки уравнений Боголюбова - Борна - Грина - Кирквуда Ивона (ББГКИ-цепочка) для $s$-частичных функций распре-





\end{document}